\documentclass[conference]{IEEEtran}
\IEEEoverridecommandlockouts

\usepackage{amsmath,amssymb}
\usepackage{graphicx}
\usepackage{booktabs}
\usepackage{xcolor}
\usepackage{cite}
\usepackage{url}

% Placeholder for numbers that must come from the confirmed benchmark run.
% Every \tbd{} is replaced only after the engine results are final.
\newcommand{\tbd}[1]{\textcolor{red}{[TBD: #1]}}

\begin{document}

\title{Learner-State-Aware Generative User Interfaces:\\
A Closed-Loop Control Architecture for Adaptive Learning Environments}

\author{\IEEEauthorblockN{Omanand Prashant Swami}
\IEEEauthorblockA{\tbd{affiliation}\\
\tbd{email}}}

\maketitle

\begin{abstract}
Intelligent tutoring systems estimate what a learner knows but present a
fixed interface, while generative user interfaces (GenUI) synthesize
interfaces at runtime without any model of the learner. We propose a
closed-loop architecture that couples the two. Bayesian Knowledge Tracing
estimates mastery of a knowledge component; a discrete-time PID governor with
anti-windup clamping and a filtered derivative converts the mastery deficit
into a structural complexity budget; a fast policy selects discrete
scaffolding affordances; and a foundation model emits a declarative JSON
component tree that must satisfy the budget, with a deterministic fallback
when schema or budget checks fail. A computable metric of interface
complexity closes the loop between the budget and the rendered tree. We
evaluate the architecture in silico on 1,000 synthetic learner trajectories
drawn from three learner archetypes and replayed through four experimental
arms. We report constraint adherence, interface jitter, recovery after error
streaks and schema compliance \tbd{headline numbers}. We further show that
standard BKT saturates after long success runs, so that a subsequent streak
of errors leaves the interface unchanged, and evaluate forgetting and a
CUSUM lapse detector as remedies. All claims concern algorithmic behavior;
no learning gains are claimed.
\end{abstract}

\begin{IEEEkeywords}
generative user interfaces, Bayesian knowledge tracing, PID control,
cognitive load, adaptive learning, simulation
\end{IEEEkeywords}

% ---------------------------------------------------------------------------
\section{Introduction}

Intelligent tutoring systems (ITS) have long modeled learners precisely.
Knowledge tracing maintains a probability that a learner has mastered each
knowledge component \cite{corbett1995knowledge,koedinger2012kli}, and
tutors built on it have been deployed at scale in mathematics
\cite{ritter2007cognitive,vanlehn2011relative}. The interface of such
systems, however, is authored in advance: the learner model selects the next
problem, not the form in which the problem is presented.

Generative user interfaces take the opposite position. A foundation model
produces a task-specific interface for each request
\cite{leviathan2026genui}, including interactive courseware
\cite{tu2026maicui}. These systems are open loop: the model does not observe
the learner's state, and nothing bounds how complex the generated interface
becomes.

Cognitive load theory explains why this matters. The same material can help
a novice and hinder an expert (the expertise reversal effect)
\cite{kalyuga2003expertise,chen2017element}, and guidance should fade as
competence grows \cite{renkl2003structuring,salden2010expertise}. An
interface that adapts to the learner therefore needs both a learner model and
a mechanism that holds the interface to a complexity appropriate to that
model, without the abrupt layout changes that make adaptive interfaces hard
to use \cite{gajos2008predictability,findlater2004menus}.

This paper makes three contributions:
\begin{enumerate}
\item A closed-loop architecture in which a BKT mastery estimate drives a PID
governor whose output is a structural complexity budget that a generative
model must satisfy, with schema validation and a deterministic fallback.
\item A computable structural complexity metric over a declarative component
tree for linear-equation tasks, built from features linked in prior work to
cognitive load and visual complexity.
\item An in-silico benchmark of 1,000 synthetic learners across four arms
and ablations, including an analysis of BKT saturation after success runs
and two remedies.
\end{enumerate}

% ---------------------------------------------------------------------------
\section{Related Work}

\subsection{Learner modeling}
Bayesian Knowledge Tracing (BKT) models each knowledge component as a
two-state hidden Markov model with prior, learning, guess and slip parameters
\cite{corbett1995knowledge}. Extensions individualize the prior
\cite{pardos2010individualization} and the learning rate
\cite{yudelson2013individualized}, add forgetting
\cite{qiu2011time,khajah2016deep}, or contextualize guess and slip
\cite{baker2008contextual}. Its identifiability and degeneracy have been
analyzed in detail \cite{beck2007identifiability,vandesande2013properties,doroudi2017misidentified}.
Logistic models such as PFA \cite{pavlik2009pfa,gong2011construct} and deep
models \cite{piech2015dkt} are alternatives; BKT remains competitive when
extended \cite{khajah2016deep} and is reviewed in
\cite{pelanek2017bkt,desmarais2012review}. We use BKT because its state is a
calibrated probability that can serve directly as a control variable.

\subsection{Adaptive and generative interfaces}
Adaptive hypermedia \cite{brusilovsky2001adaptive} and model-based interface
generation \cite{gajos2010supple} adapt presentation to users, but
adaptation has costs: unpredictable changes reduce usability
\cite{gajos2008predictability,gajos2006design,lavie2010adaptive}. Recent GenUI
systems let a large language model produce the interface itself
\cite{leviathan2026genui,tu2026maicui}. To our knowledge, none of these
systems bounds the generated interface by a learner model.

\subsection{Cognitive load and scaffolding}
Cognitive load theory \cite{sweller1988cognitive,sweller1998cognitive,sweller2019cognitive}
attributes intrinsic load to element interactivity \cite{sweller2010element}.
Worked examples reduce load for algebra novices \cite{sweller1985worked} and
should fade into problem solving \cite{renkl2003structuring,salden2010expertise};
concrete representations should fade into symbolic ones
\cite{fyfe2014concreteness}. For linear equations, the balance model supports
early understanding but breaks down with negative quantities
\cite{vlassis2002balance,otten2019balance}, and learners hold persistent
misconceptions about equality and operating on the unknown
\cite{kieran1981equality,herscovics1994gap,alibali2007equal,booth2014persistent}.

\subsection{Control and constrained generation}
PID control with anti-windup and derivative filtering is standard practice
\cite{astrom2006advanced,astrom2008feedback,astrom1989windup}, and feedback
control has been applied to computing systems \cite{hellerstein2004feedback}
and to difficulty adjustment in games \cite{hunicke2005dda}. Constrained
decoding guarantees well-formed model output
\cite{willard2023guided,scholak2021picard,geng2023grammar,poesia2022synchromesh},
though schema-constrained output is not perfectly reliable in practice
\cite{geng2025jsonschemabench} and format restrictions can affect output
quality \cite{tam2024speak}.

% ---------------------------------------------------------------------------
\section{Architecture}

The loop has four layers (Fig.~\tbd{architecture figure}): a psychometric
estimator, a complexity governor, a fast policy, and a generative plant whose
output passes a validation gate before a deterministic client renders it.

\subsection{Mastery estimation}
After response $r_n\in\{0,1\}$ the mastery estimate is updated by Bayes' rule
and the learning transition \cite{corbett1995knowledge}:
\begin{align}
P(L_n\mid r_n{=}1) &= \frac{P(L_{n})(1-P(S))}{P(L_{n})(1-P(S)) + (1-P(L_{n}))P(G)},\\
P(L_n\mid r_n{=}0) &= \frac{P(L_{n})P(S)}{P(L_{n})P(S) + (1-P(L_{n}))(1-P(G))},\\
P(L_{n+1}) &= P(L_n\mid r_n)\,(1-P(F)) + \bigl(1-P(L_n\mid r_n)\bigr)P(T),
\end{align}
where $P(F)=0$ gives standard BKT and $P(F)>0$ adds forgetting
\cite{khajah2016deep}. Default parameters are $P(L_0)=0.15$, $P(T)=0.12$,
$P(G)=0.20$, $P(S)=0.08$. They are design values, not fitted; they satisfy
the bounds $P(G)<0.3$, $P(S)<0.1$ \cite{corbett1995knowledge} and the
non-degeneracy conditions $P(G)+P(S)<1$, $P(T)<1-P(S)/(1-P(G))$
\cite{vandesande2013properties}.

\subsection{Complexity governor}
The setpoint is the conventional mastery threshold $P^*=0.95$
\cite{corbett1995knowledge,zhang2025mastery}. The tracking error is normalized
on each side of the setpoint so that deficit and surplus have equal
authority:
\begin{equation}
e_n=\begin{cases}(P^*-P(L_n))/P^* & P(L_n)\le P^*\\
(P^*-P(L_n))/(1-P^*) & P(L_n)>P^*\end{cases},\quad e_n\in[-1,1].
\end{equation}
The control effort is
\begin{equation}
u_n = K_p e_n + K_i S_n + K_d D_n,
\end{equation}
with the clamped integral $S_n=\mathrm{clamp}(S_{n-1}+e_n,\,-S_{\max},\,S_{\max})$
for anti-windup \cite{astrom1989windup} and the first-order filtered derivative
$D_n=\gamma(e_n-e_{n-1})+(1-\gamma)D_{n-1}$ \cite{astrom2006advanced}. An
inverted logistic maps effort to the complexity budget,
\begin{equation}
M_I^* = \bigl(1+e^{k u_n}\bigr)^{-1},
\end{equation}
which is quantized into four levels $\ell_n=\min(4,1+\lfloor 4M_I^*\rfloor)$. A
hysteresis deadband $\Delta_h$ commits a new budget only when it differs from
the last committed budget by at least $\Delta_h$.

\paragraph*{Gain selection}
Classical tuning rules \cite{ziegler1942optimum,cohen1953retarded,astrom2004revisiting,skogestad2003simple}
assume an identified plant response. In our benchmark the interface does not
alter the simulated learner's responses (Section~\ref{sec:method}), so there is
no plant response to identify. Instead, $K_p$, $K_i$, $K_d$, $\gamma$,
$S_{\max}$ and $\Delta_h$ were selected by minimizing an ITAE-type tracking
cost \cite{graham1953itae} plus a jitter penalty on a separate tuning cohort;
all reported results use a disjoint test cohort. Selected values:
\tbd{tuned gains from results/}. The logistic slope $k$ is chosen so that the
steady-state effort range maps to approximately $[0.05,0.95]$
\tbd{re-derive $k$ for $P^*=0.95$}.

\subsection{Fast policy}
A fast policy maps $(P(L_n), M_I^*, \ell_n, \text{recent errors})$ to discrete
affordances: scaffolding mode (worked example, guided scaffold, open
symbolic), whether to show the balance-scale manipulative, whether hints are
enabled, and an instructional density limit. The intended implementation is
an open-weight decision model \cite{laya2026}; the benchmark uses a
deterministic rule-based stand-in with the same interface, so no latency
claims about that model are made. The split between a fast decider and a
slow generator follows the System~1 / System~2 analogy
\cite{kahneman2011thinking}.

\subsection{Generative plant and validation gate}
The plant (\tbd{Gemini model and configuration}) receives the budget, the
policy decision and the equation and emits a JSON component tree. Nodes are
drawn from a closed set of six components: \textsc{ConcreteBalanceScale},
\textsc{WorkedSolutionStep}, \textsc{ScaffoldedOperationPad},
\textsc{SymbolicEquationWorkspace}, \textsc{SteppedHintAccordion} and
\textsc{CartesianVerificationPlane}. The client renders these components
directly and never executes generated code. Each emission is parsed against
the schema \cite{pezoa2016jsonschema}; if parsing fails, or if
$|M_I-M_I^*|>\varepsilon$ with $\varepsilon=0.05$, a precompiled fallback
tree nearest to the budget is rendered instead.

\subsection{Structural complexity metric}
For a tree with $\rho$ interactive elements, abstraction level $\alpha$ and
nesting depth $\delta$,
\begin{equation}
M_I = w_1\frac{\rho-1}{7} + w_2\frac{\alpha-1}{3} + w_3\frac{\delta-1}{3},
\quad \rho\in[1,8],\ \alpha,\delta\in[1,4].
\end{equation}
Each feature has a basis in prior work: element interactivity and choice
count raise load \cite{sweller2010element,hick1952rate,cowan2001magical}
and display density is a classic complexity measure \cite{tullis1983formatting};
the abstraction ladder from balance scale to operation pad to symbolic
workspace follows concreteness fading \cite{fyfe2014concreteness}; and
deeper structures cost users more \cite{kiger1984depth,larson1998web}.
Placing the Cartesian plane at the top of the ladder is our design choice.
With no data to fit weights, we use equal weights $w_i=1/3$, which are robust
in that setting \cite{dawes1979robust}, and test sensitivity to the weights
(Section~\ref{sec:results}). $M_I$ is a structural proxy; it has not been
validated against measured cognitive load.

% ---------------------------------------------------------------------------
\section{Method}
\label{sec:method}

\subsection{Synthetic learners}
Simulated learners are an established way to test instructional designs
before human studies \cite{vanlehn1994simulated,kaser2024simulated}, and
BKT-generated data have been used to evaluate mastery assessment
\cite{fancsali2013optimal,slater2018degree}. Each learner is a two-state
HMM whose parameters are drawn uniformly from archetype ranges
(Table~\ref{tab:archetypes}). The ranges for the novice and guesser
archetypes exceed the tighter bounds of \cite{corbett1995knowledge} by design
to represent noisy learners, while respecting the non-degeneracy conditions
of \cite{vandesande2013properties}. The cohort has 1,000 learners
(300/400/300) with 40 items each. Responses are sampled once per learner and
replayed unchanged through every arm (paired design). The interface therefore
does not influence the simulated responses: the benchmark measures controller
behavior, not learning.

\begin{table}[t]
\centering
\caption{Learner archetypes (uniform parameter ranges).}
\label{tab:archetypes}
\begin{tabular}{lcccc}
\toprule
Archetype & $P(L_0)$ & $P(T)$ & $P(G)$ & $P(S)$\\
\midrule
Struggling novice & .02--.08 & .03--.07 & .10--.20 & .20--.30\\
Inconsistent guesser & .05--.20 & .10--.20 & .15--.40 & .08--.30\\
Fast master & .15--.30 & .30--.40 & .15--.25 & .01--.03\\
\bottomrule
\end{tabular}
\end{table}

\subsection{Arms}
(1) \emph{Static}: a fixed level-2 interface. (2) \emph{Unconstrained GenUI}:
the plant chooses complexity from a noisy reading of mastery with no budget
check. (3) \emph{Rule-based}: step up after three consecutive correct
responses and down after two consecutive errors, an up-down rule in the
spirit of \cite{levitt1971updown,kelly2015mastery}. (4) \emph{Closed loop}:
the full architecture. Ablations remove the derivative filter, remove
anti-windup, and add a one-level slew limit.

\subsection{Plant model in the benchmark}
\tbd{If real Gemini rates are measured: describe them here. Otherwise:} The
benchmark replaces the foundation model with a surrogate that compiles the
closest admissible tree and injects faults: with probability
\tbd{0.15} it drifts off budget, and with probability \tbd{0.03} it emits
malformed output (truncated JSON, unknown component, out-of-range field,
depth overflow, injected field). These rates are assumptions, not
measurements.

\subsection{Metrics}
Constraint error $\Delta M=|M_I-M_I^*|$; schema compliance (valid trees over
attempts); interface jitter $J=\frac{1}{N-1}\sum_n|M_I(n)-M_I(n-1)|$; recovery
time after a five-error streak; and fallback rate. Jitter operationalizes the
stability concern of \cite{gajos2008predictability}; the formula is ours.

% ---------------------------------------------------------------------------
\section{Results}
\label{sec:results}

\tbd{All numbers in this section come from the confirmed final benchmark run
(commit, seed and CSV columns to be supplied by the engine).}

\begin{table}[t]
\centering
\caption{Main results on the test cohort (mean $\pm$ 95\% CI).}
\label{tab:main}
\begin{tabular}{lcccc}
\toprule
Arm & $\Delta M$ & Jitter $J$ & Schema ok & Fallback\\
\midrule
Static & \tbd{} & \tbd{} & \tbd{} & \tbd{}\\
Unconstrained GenUI & \tbd{} & \tbd{} & \tbd{} & \tbd{}\\
Rule-based & \tbd{} & \tbd{} & \tbd{} & \tbd{}\\
Closed loop & \tbd{} & \tbd{} & \tbd{} & \tbd{}\\
\bottomrule
\end{tabular}
\end{table}

\subsection{Ablations}
\tbd{derivative filter, anti-windup, slew limit}.

\subsection{Sensitivity}
\tbd{gains at $\pm 50\%$; metric weights: equal, $0.40/0.35/0.25$, and each
term dropped}.

\subsection{BKT saturation and lapse handling}
Under correct responses the BKT estimate converges to a stable fixed point at
one \cite{vandesande2013properties}. With the default parameters the estimate
exceeds $1-10^{-6}$ after ten correct responses and five subsequent errors
lower it only to about $0.98$; after roughly forty correct responses it
equals $1.0$ in double precision and no further evidence can lower it. The
interface therefore stays at its highest level through a streak of errors.
We compare two remedies: forgetting, $P(F)>0$ \cite{qiu2011time,khajah2016deep},
and a CUSUM detector \cite{page1954cusum,basseville1993detection} on the
log-likelihood ratio between ``not known'' and the tracker's prediction,
which resets the estimate when an alarm fires. \tbd{results on the
forgetting cohort}. Weighting recent responses more heavily
\cite{galyardt2015recent} is a third option we leave for future work.

% ---------------------------------------------------------------------------
\section{Discussion and Limitations}

The benchmark tests whether the controller does what it is designed to do; it
does not test whether learners benefit. Because simulated responses do not
depend on the interface, the loop is closed only through the mastery
estimate, not through the learner. $M_I$ is a structural proxy that has not
been validated against measured load, for which instruments exist
\cite{leppink2013instrument,paas1992training,ayres2006subjective}. BKT
parameters are design values rather than fits to data such as ASSISTments
\cite{badrinath2021pybkt}. Long error streaks may reflect misconceptions
rather than slips \cite{booth2014persistent}, or a learner who will not reach
mastery \cite{beck2013wheel}; the controller does not distinguish these.
\tbd{plant: surrogate vs.\ measured Gemini rates}.

\section{Conclusion}
\tbd{one paragraph after results are final}.

\bibliographystyle{IEEEtran}
\bibliography{../docs/references}

\end{document}
